% §\,4. Некоторые дополнения
% Зорич, Математический анализ, Том I
% Книжные страницы 23--31 (PDF-страницы 35--43)

\section*{\S\,4. Некоторые дополнения}

\subsection*{1. Мощность множества (кардинальные числа)}

Говорят, что множество $X$ \textit{равномощно} множеству $Y$, если существует
биективное отображение $X$ на $Y$, т.\,е. каждому элементу $x \in X$ сопоставляется
элемент $y \in Y$, причем различным элементам множества $X$ отвечают различные
элементы множества $Y$ и каждый элемент $y \in Y$ сопоставлен некоторому элементу
множества $X$.

Описательно говоря, каждый элемент $x \in X$ сидит на своем месте $y \in Y$,
все элементы $X$ сидят и свободных мест $y \in Y$ нет.

Ясно, что введенное отношение $X \mathscr{R} Y$ является отношением \textit{эквивалентности}, поэтому мы будем, как и договаривались, писать в этом случае $X \sim Y$
вместо $X \mathscr{R} Y$.

Отношение равномощности разбивает совокупность всех множеств на
классы эквивалентных между собой множеств. Множества одного класса
эквивалентности имеют одинаковое количество элементов (равномощны), а
разных --- разное.

Класс, которому принадлежит множество $X$, называется \textit{мощностью множества} $X$, а также \textit{кардиналом} или \textit{кардинальным числом множества} $X$ и
обозначается символом $\operatorname{card} X$. Если $X \sim Y$, то пишут $\operatorname{card} X = \operatorname{card} Y$.

Смысл этой конструкции в том, что она позволяет сравнивать количества элементов множеств, не прибегая к промежуточному счету, т.\,е. к
измерению количества путем сравнения с натуральным рядом чисел $\mathbb{N} =
{} = \{1, 2, 3, \ldots\}$. Последнее, как мы вскоре увидим, иногда принципиально
невозможно.

Говорят, что кардинальное число множества $X$ \textit{не больше} кардинального
числа множества $Y$, и пишут $\operatorname{card} X \leqslant \operatorname{card} Y$, если $X$ равномощно некоторому подмножеству множества $Y$.

Итак,
\[
  (\operatorname{card} X \leqslant \operatorname{card} Y) := (\exists Z \subset Y \mid \operatorname{card} X = \operatorname{card} Z).
\]

Если $X \subset Y$, то ясно, что $\operatorname{card} X \leqslant \operatorname{card} Y$. Однако, оказывается, соотношение $X \subset Y$ не мешает неравенству $\operatorname{card} Y \leqslant \operatorname{card} X$, даже если $X$ есть собственное подмножество $Y$.

Например, соответствие $x \mapsto \dfrac{x}{1 - |x|}$ есть биективное отображение
промежутка $-1 < x < 1$ числовой оси $\mathbb{R}$ на всю эту ось.

Возможность для множества быть равномощным своей части является
характерным признаком бесконечных множеств, который Дедекинд\footnote{Р.\,Дедекинд (1831--1916) --- немецкий математик-алгебраист, принявший активное
участие в развитии теории действительного числа. Впервые предложил аксиоматику
множества натуральных чисел, называемую обычно аксиоматикой Пеано --- по имени Дж.\;Пеано (1858--1932), итальянского математика, сформулировавшего ее несколько позже.} даже
предложил считать определением бесконечного множества. Таким образом,
множество называется \textit{конечным} (по Дедекинду), если оно не равномощно
никакому своему собственному подмножеству; в противном случае оно
называется \textit{бесконечным}.

Подобно тому, как отношение неравенства упорядочивает действительные числа на числовой прямой, введенное отношение неравенства упорядочивает мощности или кардинальные числа множеств. А именно, можно
доказать, что справедливы следующие свойства построенного отношения:

$1^\circ$ \; $(\operatorname{card} X \leqslant \operatorname{card} Y) \land (\operatorname{card} Y \leqslant \operatorname{card} Z) \Rightarrow (\operatorname{card} X \leqslant \operatorname{card} Z)$ \; (очевидно);

$2^\circ$ \; $(\operatorname{card} X \leqslant \operatorname{card} Y) \land (\operatorname{card} Y \leqslant \operatorname{card} X) \Rightarrow (\operatorname{card} X = \operatorname{card} Y)$ \; (теорема Шрёдера--Берн\-штей\-на\footnote{Ф.\,Бернштейн (1878--1956) --- немецкий математик, ученик Г.\;Кантора; Э.\;Шрёдер
(1841--1902) --- немецкий математик.});

$3^\circ$ \; $\forall X\; \forall Y$ \; $(\operatorname{card} X \leqslant \operatorname{card} Y) \lor (\operatorname{card} Y \leqslant \operatorname{card} X)$ \; (теорема Кантора).

\medskip

Таким образом, класс кардинальных чисел оказывается линейно упорядоченным.

Говорят, что мощность множества $X$ \textit{меньше} мощности множества $Y$, и
пишут $\operatorname{card} X < \operatorname{card} Y$, если $\operatorname{card} X \leqslant \operatorname{card} Y$ и в то же время $\operatorname{card} X \neq \operatorname{card} Y$.

Итак, $(\operatorname{card} X < \operatorname{card} Y) := (\operatorname{card} X \leqslant \operatorname{card} Y) \land (\operatorname{card} X \neq \operatorname{card} Y)$.

Пусть, как и прежде, $\varnothing$ --- знак пустого множества, а $\mathscr{P}(X)$ --- символ
множества всех подмножеств множества $X$. Имеет место следующая
открытая Кантором

\textsc{Теорема.} \; $\operatorname{card} X < \operatorname{card} \mathscr{P}(X)$.

$\blacktriangleleft$ \; Для пустого множества $\varnothing$ утверждение очевидно, поэтому в дальнейшем можно считать, что $X \neq \varnothing$.

Поскольку $\mathscr{P}(X)$ содержит все одноэлементные подмножества $X$,
$\operatorname{card} X \leqslant {} \leqslant \operatorname{card} \mathscr{P}(X)$.

Для доказательства теоремы теперь достаточно установить, что $\operatorname{card} X \neq {} \neq \operatorname{card} \mathscr{P}(X)$, если $X \neq \varnothing$.

Пусть, вопреки утверждению, существует биективное отображение
$f \colon X \to \mathscr{P}(X)$. Рассмотрим множество $A = \{x \in X \mid x \notin f(x)\}$ тех элементов
$x \in X$, которые не содержатся в сопоставленном им множестве $f(x) \in \mathscr{P}(X)$.
Поскольку $A \in \mathscr{P}(X)$, то найдется элемент $a \in X$ такой, что $f(a) = A$. Для
элемента $a \in X$ невозможно ни соотношение $a \in A$ (по определению $A$),
ни соотношение $a \notin A$ (опять-таки по определению $A$). Мы вступаем в
противоречие с законом исключенного третьего. $\blacktriangleright$

Эта теорема, в частности, показывает, что если бесконечные множества
существуют, то и <<бесконечности>> бывают разные.

\subsection*{2. Об аксиоматике теории множеств}

Цель настоящего пункта --- дать интересующемуся читателю представление о системе аксиом, описывающих свойства математического объекта, называемого \textit{множеством}, и продемонстрировать простейшие
следствия этих аксиом.

$1^\circ$ \textsc{Аксиома объемности}. \textit{Множества $A$ и $B$ равны тогда и только тогда, когда они имеют одни и те же элементы.}

Это означает, что мы отвлекаемся от всех прочих свойств объекта <<множество>>,
кроме свойства иметь данные элементы. На практике это означает, что если мы
желаем установить, что $A = B$, то мы должны проверить, что $\forall x\, ((x \in A) \Leftrightarrow (x \in B))$.

$2^\circ$ \textsc{Аксиома выделения}. \textit{Любому множеству $A$ и свойству $P$ отвечает множество $B$, элементы которого суть те и только те элементы множества $A$, которые
обладают свойством~$P$.}

Короче, утверждается, что если $A$ --- множество, то и $B = \{x \in A \mid P(x)\}$ --- тоже
множество.

Эта аксиома очень часто используется в математических конструкциях, когда мы
выделяем из множеств подмножества, состоящие из элементов, обладающих тем или
иным свойством.

Например, из аксиомы выделения следует, что существует пустое подмножество
$\varnothing_X = \{x \in X \mid x \neq x\}$ в любом множестве $X$, а с учетом аксиомы объемности
заключаем, что для любых множеств $X$ и $Y$ выполнено $\varnothing_X = \varnothing_Y$, т.\,е. пустое
множество единственно. Его обозначают символом $\varnothing$.

Из аксиомы выделения следует также, что если $A$ и $B$ --- множества, то
$A \setminus B = \{x \mid x \notin B\}$ --- тоже множество. В частности, если $M$ --- множество и $A$ --- его подмножество, то $C_M A$ --- тоже множество.

$3^\circ$ \textsc{Аксиома объединения}. \textit{Для любого множества $M$ множеств существует множество $\bigcup M$, называемое объединением множества $M$, состоящее из тех и только
тех элементов, которые содержатся в элементах множества~$M$.}

Если вместо слов <<множество множеств>> сказать <<семейство множеств>>, то
аксиома объединения приобретает несколько более привычное звучание: существует
множество, состоящее из элементов множеств семейства. Таким образом, объединение множества есть множество, причем $x \in \bigcup M \Leftrightarrow \exists X\; ((X \in M) \land (x \in X))$.

Аксиома объединения с учетом аксиомы выделения позволяет определить
\textit{пересечение множества $M$} (семейства множеств) как множество
\[
  \bigcap M := \{x \in \bigcup M \mid \forall X\; ((X \in M) \Rightarrow (x \in X))\}.
\]

$4^\circ$ \textsc{Аксиома пары}. \textit{Для любых множеств $X$ и $Y$ существует множество $Z$ такое,
что $X$ и $Y$ являются его единственными элементами.}

Множество $Z$ обозначается через $\{X, Y\}$ и называется \textit{неупорядоченной парой} множеств $X$ и $Y$. Множество $Z$ состоит из одного элемента, если $X = Y$.

Как мы уже отмечали, \textit{упорядоченная} пара $(X, Y)$ множеств отличается от
неупорядоченной наличием какого-либо признака у одного из множеств пары. Например,
$(X, Y) := \{\{X, X\}, \{X, Y\}\}$.

Итак, неупорядоченная пара позволяет ввести упорядоченную пару, а
упорядоченная пара позволяет ввести прямое произведение множеств, если
воспользоваться аксиомой выделения и следующей важной аксиомой.

$5^\circ$ \textsc{Аксиома множества подмножеств}. \textit{Для любого множества $X$ существует множество $\mathscr{P}(X)$, состоящее из тех и только тех элементов, которые являются
подмножествами множества~$X$.}

Короче говоря, существует множество всех подмножеств данного множества.

Теперь можно проверить, что упорядоченные пары $(x, y)$, где $x \in X$, а $y \in Y$,
действительно образуют множество
\[
  X \times Y := \{p \in \mathscr{P}(\mathscr{P}(X) \cup \mathscr{P}(Y)) \mid p = (x, y) \land (x \in X) \land (y \in Y)\}.
\]

Аксиомы $1^\circ$--$5^\circ$ ограничивают возможность формирования новых множеств. Так,
в множестве $\mathscr{P}(X)$ по теореме Кантора (о том, что $\operatorname{card} X < \operatorname{card} \mathscr{P}(X)$) имеется
элемент, не принадлежащий $X$, поэтому <<множества>> \textit{всех множеств} не существует.
А ведь именно на этом <<множестве>> держится парадокс Рассела.

Для того чтобы сформулировать следующую аксиому, введем понятие
\textit{последователя} $X^+$ множества $X$. Положим по определению $X^+ = X \cup \{X\}$. Короче,
к $X$ добавлено одноэлементное множество $\{X\}$.

Далее, множество назовем \textit{индуктивным}, если оно содержит в качестве элемента
пустое множество и последователь любого своего элемента.

$6^\circ$ \textsc{Аксиома бесконечности}. \textit{Индуктивные множества существуют.}

Аксиома бесконечности позволяет с учетом аксиом $1^\circ$--$4^\circ$ создать эталонную
модель множества $\mathbb{N}_0$ натуральных чисел (по фон Нейману\footnote{Дж.\;фон Нейман (1903--1957) --- американский математик. Работы по функциональному анализу, математическим основаниям квантовой механики, топологическим группам, теории игр, математической логике. Руководил созданием первых ЭВМ.}), определив $\mathbb{N}_0$ как
пересечение индуктивных множеств, т.\,е. как наименьшее индуктивное множество.
Элементами $\mathbb{N}_0$ являются множества
\[
  \varnothing, \quad \varnothing^+ = \varnothing \cup \{\varnothing\} = \{\varnothing\}, \quad
  \{\varnothing\}^+ = \{\varnothing\} \cup \{\{\varnothing\}\}, \quad \ldots,
\]
которые и являются моделью того, что мы обозначаем символами $0, 1, 2, \ldots$ и
называем натуральными числами.

$7^\circ$ \textsc{Аксиома подстановки}. \textit{Пусть $\mathscr{F}(x, y)$ --- такое высказывание (точнее,
формула), что при любом $x_0$ из множества $X$ существует и притом единственный
объект $y_0$ такой, что $\mathscr{F}(x_0, y_0)$ истинно. Тогда объекты~$y$, для каждого из которых
существует элемент $x \in X$ такой, что $\mathscr{F}(x, y)$ истинно, образуют множество.}

Этой аксиомой при построении анализа мы пользоваться не будем.

Аксиомы $1^\circ$--$7^\circ$ составляют аксиоматику теории множеств, известную как аксиоматика Цермело--Френкеля\footnote{Э.\,Цермело (1871--1953) --- немецкий математик; А.\,Френкель (1891--1965) --- немецкий, затем израильский математик.}.

К ней обычно добавляется еще одна, независимая от аксиом $1^\circ$--$7^\circ$ и часто используемая в анализе

$8^\circ$ \textsc{Аксиома выбора}. \textit{Для любого семейства непустых попарно непересекающихся
множеств существует множество $C$ такое, что, каково бы ни было множество $X$
данного семейства, множество $X \cap C$ состоит из одного элемента.}

Иными словами, из каждого множества семейства можно выбрать в точности по
одному представителю так, что выбранные элементы составят множество $C$.

Аксиома выбора, известная в математике как аксиома Цермело, вызвала горячие
дискуссии специалистов.

\subsection*{3. Замечания о структуре математических высказываний и записи
их на языке теории множеств}

В языке теории множеств имеются два базисных или, как говорят, атомарных типа математических высказываний:
утверждение $x \in A$ о том, что объект $x$ есть элемент множества $A$, и утверждение $A = B$ о том, что множества $A$ и $B$ совпадают. (Впрочем, с учетом
аксиомы объемности второе утверждение является комбинацией утверждений первого типа: $(x \in A) \Leftrightarrow (x \in B)$.)

Сложное высказывание или сложная логическая формула строятся из
атомарных посредством логических операторов --- связок $\neg$, $\land$, $\lor$, $\Rightarrow$ и
кванторов $\forall$, $\exists$ с использованием скобок ( ). При этом формирование сколь
угодно сложного высказывания и его записи сводится к выполнению
следующих элементарных логических операций:

a) образование нового высказывания путем постановки отрицания перед
некоторым высказыванием и заключение результата в скобки;

b) образование нового высказывания путем постановки необходимой
связки $\land$, $\lor$, $\Rightarrow$ между двумя высказываниями и заключение результата в
скобки;

c) образование высказывания <<для любого объекта $x$ выполнено свойство $P$>> (что записывают в виде $\forall x\, P(x)$) или высказывания <<найдется объект $x$, обладающий свойством $P$>> (что записывают в виде $\exists x\, P(x)$).

\medskip

Например, громоздкая запись
\[
  \exists x\; (P(x) \land (\forall y\; ((P(y)) \Rightarrow (y = x))))
\]
означает, что найдется объект $x$, обладающий свойством $P$ и такой, что если
$y$ --- любой объект, обладающий свойством $P$, то $y = x$. Короче: существует
и притом единственный объект $x$, обладающий свойством~$P$. Обычно это
высказывание обозначают в виде $\exists!\, x\, P(x)$, и мы будем использовать такое
сокращение.

Для упрощения записи высказывания, как уже отмечалось, стараются
опустить столько скобок, сколько это возможно без потери однозначного
толкования записи. С этой целью кроме указанного ранее приоритета
операторов $\neg$, $\land$, $\lor$, $\Rightarrow$ считают, что наиболее жестко символы в формуле
связываются знаками $\in$, $=$, затем $\exists$, $\forall$ и потом связками $\neg$, $\land$, $\lor$, $\Rightarrow$.

С учетом такого соглашения теперь можно было бы написать
\[
  \exists!\, x\, P(x) := \exists x\; (P(x) \land \forall y\; (P(y) \Rightarrow y = x)).
\]

Условимся также о следующих широко используемых сокращениях:
\begin{align*}
  (\forall x \in X)\; P &:= \forall x\; (x \in X \Rightarrow P(x)), \\
  (\exists x \in X)\; P &:= \exists x\; (x \in X \land P(x)), \\
  (\forall x > a)\; P &:= \forall x\; (x \in \mathbb{R} \land x > a \Rightarrow P(x)), \\
  (\exists x > a)\; P &:= \exists x\; (x \in \mathbb{R} \land x > a \land P(x)).
\end{align*}

Здесь $\mathbb{R}$, как всегда, есть символ множества действительных чисел.

С учетом этих сокращений и правил~a), b), c) построения сложного
высказывания, например, можно будет дать однозначно трактуемую запись
\[
  \bigl(\lim_{x \to a} f(x) = A\bigr) := \forall \varepsilon > 0\; \exists \delta > 0\;
  \forall x \in \mathbb{R}\; (0 < |x - a| < \delta \Rightarrow |f(x) - A| < \varepsilon)
\]
того, что число $A$ является пределом функции $f \colon \mathbb{R} \to \mathbb{R}$ в точке $a \in \mathbb{R}$.

Быть может, наиболее важным из всего сказанного в этом параграфе
являются для нас следующие правила построения отрицания к высказыванию,
содержащему кванторы.

Отрицание к высказыванию <<для некоторого $x$ истинно $P(x)$>> означает,
что <<для любого $x$ неверно $P(x)$>>, а отрицание к высказыванию <<для любого
$x$ истинно $P(x)$>> означает, что <<найдется $x$, что неверно $P(x)$>>.

Итак,
\begin{align*}
  \neg \exists x\, P(x) &\Leftrightarrow \forall x\, \neg P(x), \\
  \neg \forall x\, P(x) &\Leftrightarrow \exists x\, \neg P(x).
\end{align*}

Напомним также (см.\ упражнения к \S\,1), что
\begin{align*}
  \neg (P \land Q) &\Leftrightarrow \neg P \lor \neg Q, \\
  \neg (P \lor Q) &\Leftrightarrow \neg P \land \neg Q, \\
  \neg (P \Rightarrow Q) &\Leftrightarrow P \land \neg Q.
\end{align*}

На основании сказанного можно заключить, что, например,
\[
  \neg ((\forall x > a)\, P) \Leftrightarrow (\exists x > a)\, \neg P.
\]

Написать в правой части последнего соотношения $(\exists x \leqslant a)\, \neg P$ было бы,
конечно, ошибочно.

В самом деле,
\begin{align*}
  \neg ((\forall x > a)\, P) &:= \neg (\forall x\; (x \in \mathbb{R} \land x > a \Rightarrow P(x))) \Leftrightarrow \\
  &\Leftrightarrow \exists x\; \neg (x \in \mathbb{R} \land x > a \Rightarrow P(x)) \Leftrightarrow \\
  &\Leftrightarrow \exists x\; ((x \in \mathbb{R} \land x > a) \land \neg P(x)) =: (\exists x > a)\, \neg P.
\end{align*}

Если учесть указанную выше структуру произвольного высказывания, то
теперь с использованием построенных отрицаний простейших высказываний
можно было бы построить отрицание любого конкретного высказывания.

Например,
\[
  \neg \bigl(\lim_{x \to a} f(x) = A\bigr) \Leftrightarrow
  \exists \varepsilon > 0\; \forall \delta > 0\; \exists x \in \mathbb{R}\;
  (0 < |x - a| < \delta \land |f(x) - A| \geqslant \varepsilon).
\]

Практическая важность правильного построения отрицания связана, в
частности, с методом доказательства от противного, когда истинность
некоторого утверждения $P$ извлекают из того, что утверждение $\neg P$ ложно.

\bigskip

\textbf{Упражнения}

\medskip

\textbf{1.}\; a) Установите равномощность отрезка $\{x \in \mathbb{R} \mid 0 \leqslant x \leqslant 1\}$ и интервала
$\{x \in \mathbb{R} \mid 0 < x < 1\}$ числовой прямой $\mathbb{R}$ как с помощью теоремы Шрёдера--Берн\-штей\-на,
так и непосредственным предъявлением нужной биекции.

b) Разберите следующее доказательство теоремы Шрёдера--Берн\-штей\-на:
\[
  (\operatorname{card} X \leqslant \operatorname{card} Y) \land (\operatorname{card} Y \leqslant \operatorname{card} X) \Rightarrow
  (\operatorname{card} X = \operatorname{card} Y).
\]

$\blacktriangleleft$ \; Достаточно доказать, что если множества $X$, $Y$, $Z$ таковы, что $X \supset Y \supset Z$ и
$\operatorname{card} X = \operatorname{card} Z$, то $\operatorname{card} X = \operatorname{card} Y$. Пусть $f \colon X \to Z$ --- биективное отображение. Тогда биекция $g \colon X \to Y$ может быть задана, например, следующим образом:
\[
  g(x) = \begin{cases}
    f(x), & \text{если } x \in f^n(X) \setminus f^n(Y) \text{ для некоторого } n \in \mathbb{N}, \\
    x & \text{в противном случае.}
  \end{cases}
\]
Здесь $f^n = f \circ \ldots \circ f$ --- $n$-я итерация отображения $f$, а $\mathbb{N}$ --- множество натуральных
чисел. $\blacktriangleright$

\textbf{2.}\; a) Исходя из определения пары, проверьте, что данное в пункте~2 определение
прямого произведения $X \times Y$ множеств $X$, $Y$ корректно, т.\,е. множество
$\mathscr{P}(\mathscr{P}(X) \cup \mathscr{P}(Y))$ содержит все упорядоченные пары $(x, y)$, в которых $x \in X$ и $y \in Y$.

b) Покажите, что все возможные отображения $f \colon X \to Y$ одного фиксированного множества $X$ в другое фиксированное множество $Y$ сами образуют множество
$M(X, Y)$.

\textbf{3.}\; a) Используя аксиомы объемности, пары, выделения, объединения и
бесконечности, проверьте, что для элементов множества $\mathbb{N}_0$ натуральных чисел
по фон Нейману справедливы следующие утверждения:

\begin{minipage}[t]{0.45\linewidth}
  $1^\circ$ \; $x = y \Rightarrow x^+ = y^+$;

  $3^\circ$ \; $x^+ = y^+ \Rightarrow x = y$;
\end{minipage}
\hfill
\begin{minipage}[t]{0.45\linewidth}
  $2^\circ$ \; $(\forall x \in \mathbb{N}_0)\; (x^+ \neq \varnothing)$;

  $4^\circ$ \; $(\forall x \in \mathbb{N}_0)\; (x \neq \varnothing \Rightarrow (\exists y \in \mathbb{N}_0)\; (x = y^+))$.
\end{minipage}

\medskip

b) Используя то, что $\mathbb{N}_0$ --- индуктивное множество, покажите, что для любых его
элементов $x$ и $y$ (а они в свою очередь являются множествами) справедливы
следующие соотношения:

\begin{minipage}[t]{0.65\linewidth}
  $1^\circ$ \; $\operatorname{card} x \leqslant \operatorname{card} x^+$;

  $3^\circ$ \; $\operatorname{card} x < \operatorname{card} y \Leftrightarrow \operatorname{card} x^+ < \operatorname{card} y^+$;

  $5^\circ$ \; $\operatorname{card} x < \operatorname{card} y \Rightarrow \operatorname{card} x^+ \leqslant \operatorname{card} y$;

  $7^\circ$ \; $(x \subset y) \lor (x \supset y)$.
\end{minipage}
\hfill
\begin{minipage}[t]{0.35\linewidth}
  $2^\circ$ \; $\operatorname{card} \varnothing < \operatorname{card} x^+$;

  $4^\circ$ \; $\operatorname{card} x < \operatorname{card} x^+$;

  $6^\circ$ \; $x = y \Leftrightarrow \operatorname{card} x = \operatorname{card} y$;
\end{minipage}

\medskip

c) Покажите, что в любом подмножестве $X$ множества $\mathbb{N}_0$ найдется такой
(наименьший) элемент $x_m$, что $(\forall x \in X)\; (\operatorname{card} x_m \leqslant \operatorname{card} x)$. (В случае
затруднений к этой задаче можно вернуться после прочтения главы~II.)

\textbf{4.}\; Мы будем иметь дело только с множествами. Поскольку множество,
состоящее из различных элементов, само может быть элементом другого множества,
логики все множества обычно обозначают строчными буквами. В настоящей задаче
это очень удобно.

a) Проверьте, что запись
\[
  \forall x\; \exists y\; \forall z\; (z \in y \Leftrightarrow \exists w\; (z \in w \land w \in x))
\]
выражает аксиому объединения, по которой существует множество $y$ --- объединение
множества $x$.

b) Укажите, какие аксиомы теории множеств представлены записями
\begin{gather*}
  \forall x\; \forall y\; \forall z\; ((z \in x \Leftrightarrow z \in y) \Leftrightarrow x = y), \\
  \forall x\; \forall y\; \exists z\; \forall v\; (v \in z \Leftrightarrow (v = x \lor v = y)), \\
  \forall x\; \exists y\; \forall z\; (z \in y \Leftrightarrow \forall u\; (u \in z \Rightarrow u \in x)),
\end{gather*}
\begin{multline*}
  \exists x\; (\forall y\; (\neg \exists z\; (z \in y) \Rightarrow y \in x) \land {} \\
  {} \land \forall w\; (w \in x \Rightarrow \forall u\; (\forall v\; (v \in u \Leftrightarrow (v = w \lor v = w)) \Rightarrow u \in x))).
\end{multline*}

c) Проверьте, что формула
\[
  \forall z\; (z \in f \Rightarrow (\exists x_1\; \exists y_1\; (x_1 \in x \land y_1 \in y \land z = (x_1, y_1)))) \land {}
\]
\[
  {} \land \forall x_1\; (x_1 \in x \Rightarrow \exists y_1\; \exists z\; (y_1 \in y \land z = (x_1, y_1) \land z \in f)) \land {}
\]
\[
  {} \land \forall x_1\; \forall y_1\; \forall y_2\;
  (\exists z_1\; \exists z_2\; (z_1 \in f \land z_2 \in f \land z_1 = (x_1, y_1) \land z_2 = (x_1, y_2))
  \Rightarrow y_1 = y_2)
\]
последовательно накладывает на множество $f$ три ограничения: $f$ есть подмножество $x \times y$; проекция $f$ на $x$ совпадает с $x$; каждому элементу $x_1$ из $x$ отвечает ровно
один элемент $y_1$ из $y$ такой, что $(x_1, y_1) \in f$.

Таким образом, перед нами определение отображения $f \colon x \to y$.

Этот пример еще раз показывает, что формальная запись высказывания отнюдь
не всегда бывает более короткой и прозрачной в сравнении с его записью на разговорном языке. Учитывая это обстоятельство, мы будем в дальнейшем использовать
логическую символику лишь в той мере, в какой она будет нам представляться
полезной для достижения большей компактности или ясности изложения.

\textbf{5.}\; Пусть $f \colon X \to Y$ --- отображение. Запишите логическое отрицание каждого из
следующих высказываний:

\quad a) $f$ сюръективно; \quad b) $f$ инъективно; \quad c) $f$ биективно.

\textbf{6.}\; Пусть $X$ и $Y$ --- множества и $f \subset X \times Y$. Запишите, что значит, что множество $f$
не является функцией.
